% ============================================================
%  Appendix
% ============================================================
\appendix

% ────────────────────────────────────────────────────────────
\section{LLM Prompt Templates}
\label{app:prompts}

All LLM calls in KBDebugger use the prompt templates below, loaded
verbatim from \texttt{src/kbdebugger/prompts/}.  Template variables
(e.g.\ \texttt{\$examples\_json}, \texttt{\$chunks\_json}) are
substituted at call time with serialised JSON.  Every call uses
temperature~0.0 and requests strict JSON output; no chain-of-thought
scratchpad is exposed outside the structured fields.

% ----------------------------------------------------------
\subsection{Stage 2c — Atomic Quality Decomposition (batch)}
\label{app:prompt-decompose-batch}

Used by Stage~2c to decompose up to \texttt{\$max\_qualities\_per\_chunk}
atomic quality sentences per chunk in a single batched LLM call.

\begin{lstlisting}[language={},breaklines=true,caption={},numbers=none,xleftmargin=0pt]
You are a careful assistant. Your task is to extract clear, atomic
"qualities" from EACH provided text chunk.

You will be given a JSON object containing multiple chunks. Each chunk has:
- an integer "id"
- a string "text"

You MUST return a single strict JSON object with this exact structure:
{
  "results": [
    {"id": 0, "qualities": ["...", "..."]},
    {"id": 1, "qualities": ["...", ...]}
  ]
}

Hard requirements:
- Output MUST be strict JSON (double quotes only). No markdown. No extra keys.
- The "results" array MUST contain exactly one entry per input chunk id.
- Each entry MUST preserve the original integer "id".
- Each "qualities" value MUST be a JSON array of strings (even if empty).
- Use ONLY information explicitly present in the corresponding chunk text.
  Never add external facts.
- If nothing meaningful can be extracted from a chunk, return an empty list.

What is a "quality":
- A quality is a short statement expressing exactly ONE fact (one atomic idea).
- Qualities will later be converted into S-P-O triplets. Therefore, each
  quality MUST be phrased so that a triplet can be extracted unambiguously.

Guidelines for qualities:
- Use explicit entities instead of pronouns.
  BAD: "He discovered relativity."
  GOOD: "Einstein discovered relativity."
- Preserve conditional meaning. Do NOT split a conditional into separate
  facts. Keep the full conditional as ONE quality.
  Example: "There is no entitlement to inspection if the examination
  was graded 1.0."
- Start with the most important, high-level qualities for each chunk.
- Ignore meta-information (headings, boilerplate, citations, references).
- Avoid vague phrases ("this thing", "it", "that process").
- Avoid quotation marks inside qualities.
- Return at most $max_qualities_per_chunk qualities per chunk.

Examples: $examples_json

Now extract qualities for each chunk in this JSON input:
$chunks_json
\end{lstlisting}

% ----------------------------------------------------------
\subsection{Stage 4 — LLM Novelty Classification (batch)}
\label{app:prompt-novelty-batch}

Used by Stage~4 to classify each quality sentence as
\textsc{existing}, \textsc{partially\_new}, or \textsc{new} relative
to its retrieved KG neighbours, in a single batched call.

\begin{lstlisting}[language={},breaklines=true,caption={},numbers=none,xleftmargin=0pt]
You are a careful knowledge-graph curator.

Your task: classify EACH candidate QUALITY sentence as one of:
- EXISTING: conveys no meaningful new semantic content compared to the
  neighbors.
- PARTIALLY_NEW: overlaps strongly with neighbors but adds at least one
  meaningful new detail, qualifier, scope, condition, or additional
  concept that would improve the Knowledge Graph.
- NEW: mostly not covered by the neighbors; introduces a new claim,
  relation, or aspect not present in the neighbors.

You will be given a JSON object with:
{
  "items": [
    {
      "id": <int>,
      "quality": <string>,
      "neighbors": [{"score": <float>, "sentence": <string>}, ...]
    },
    ...
  ]
}

Important guidelines:
1. Do NOT require exact lexical match. Judge semantic equivalence.
2. Prefer PARTIALLY_NEW when the QUALITY contains extra specifics not
   in the best neighbor: extra concept/entity, extra condition/trigger,
   extra scope, extra modality/strength, or extra relation type.
3. Prefer EXISTING when QUALITY is effectively a rephrase of the
   closest neighbor and adds no new constraints or concepts.
4. Use similarity scores as a weak signal:
   - max_score >= 0.90 and meaning already captured -> EXISTING likely
   - max_score medium/high but QUALITY adds a new detail -> PARTIALLY_NEW
   - max_score low/medium and content not covered -> NEW
   Scores are not decisive; content is decisive.

Output MUST be valid JSON:
{
  "results": [
    {
      "id": <int>,
      "decision": "EXISTING" | "PARTIALLY_NEW" | "NEW",
      "rationale": "string",
      "novel_spans": ["string", ...],
      "matched_neighbor_sentence": "string | null",
      "confidence": 0.0-1.0
    },
    ...
  ]
}

Hard requirements:
- Output strict JSON only. No markdown.
- Exactly one result entry per input item id.
- "novel_spans": phrases from QUALITY representing new semantic
  contribution (empty for EXISTING).
- "matched_neighbor_sentence": closest neighbor in meaning, or null.

Few-shot examples: $examples_json

Now classify these items:
$items_json
\end{lstlisting}

% ----------------------------------------------------------
\subsection{Stage 5 — Controlled-Vocabulary Triplet Extraction (batch)}
\label{app:prompt-triplets-batch}

Used by Stage~5 to extract \textsc{(Subject, Object, Predicate)}
triplets with deontic modality from a batch of quality sentences,
constrained to the declared predicate vocabulary.

\begin{lstlisting}[language={},breaklines=true,caption={},numbers=none,xleftmargin=0pt]
You are an information extraction system. Extract meaningful semantic
relationships from each sentence so they can be reviewed and inserted
into a Knowledge Graph.

Each relationship MUST be a triplet in the exact order:
(Subject, Object, Predicate)

The Predicate value is ontology-controlled. You MUST use only exact
values from this allowed predicate list: $predicates_json

Rules for Predicate selection:
- Copy the Predicate string exactly. Do not lowercase, add spaces, or
  invent variants.
- If no allowed Predicate accurately represents a sentence, return an
  empty triplets array and provide a short skipped_reason.
- Do not force-fit a sentence into a Predicate. Do not hallucinate.
- If multiple Predicates could fit, prefer the earliest in the list.

Normative language (standards documents):
- Prefer normative Predicates for obligations, recommendations,
  permissions, prohibitions, or definitions:
  Requires, Recommends, Permits, Prohibits, Defines.
- Classify each sentence's deontic modality:
  "shall" / "must"         -> "MANDATORY"
  "should"                 -> "RECOMMENDED"
  "may" / "can"            -> "OPTIONAL"
  "shall not" / "must not" -> "PROHIBITED"
  no normative wording     -> "NONE"
- For compound normative sentences, emit a SEPARATE triplet per
  distinct claim. Report the STRONGEST modality for the sentence-level
  field (priority: PROHIBITED > MANDATORY > RECOMMENDED > OPTIONAL).

Standard schema guidance:
- Subject and Object MUST come from the source sentence. Do not replace
  source terms with schema canonical node names.
- Use schema_templates and predicate_hints only to choose a Predicate
  and role pattern, not to rename entities.

You will receive a JSON array of objects, each with:
  - "id": integer
  - "sentence": string
  - "standard_schema": optional object with extracted_entities,
    inferred_node_types, predicate_hints, schema_templates,
    and grounding_confidence

You MUST respond with:
{
  "triplets_batch": [
    {
      "id": <int>,
      "sentence": <string>,
      "triplets": [[<Subject>, <Object>, <Predicate>], ...],
      "modality": <"MANDATORY"|"RECOMMENDED"|"OPTIONAL"|"PROHIBITED"|"NONE">,
      "skipped_reason": <string | null>
    },
    ...
  ]
}

Hard requirements:
- Output strict JSON only. No markdown, no prose outside JSON.
- Exactly one entry per input id.
- Each triplet: exactly three strings in order Subject, Object, Predicate.
- Every Predicate must be from the allowed list.
- If no valid triplets: "triplets": [] and set skipped_reason.

Normalization (Subject and Object):
- Remove leading articles ("a", "an", "the").
- Prefer singular nouns when meaning is preserved.
- Short, specific noun phrases; normal casing for common nouns,
  preserved capitalisation for proper names and acronyms.
- Never use "There", "It", "This", or "That" as Subject or Object.

Examples: $examples_json

Now process this input:
$payload_json
\end{lstlisting}

% ----------------------------------------------------------
\subsection{Stage 6 — Conflict \& Tension Judge}
\label{app:prompt-conflict-judge}

Used by the comparison engine to adjudicate candidate inter-standard
conflicts, classifying each as \textsc{agree}, \textsc{unrelated},
\textsc{tension}, or \textsc{contradict}.

\begin{lstlisting}[language={},breaklines=true,caption={},numbers=none,xleftmargin=0pt]
You are an expert reviewer of AI standards and guidelines (ISO/IEC,
Fraunhofer, EU AI Act and similar). You compare pairs of statements
extracted from different documents and judge whether they agree or
pull against each other.

You will receive a JSON array of candidate pairs. Each candidate has:
  - "id": string identifier (copy it back unchanged)
  - "type": structural signal (MODALITY_CONFLICT, DEFINITION_DIVERGENCE,
    TAXONOMY_CONFLICT, VALUE_CONFLICT)
  - "side_a": {"doc": document name, "text": verbatim statement,
    "modality": optional obligation strength}
  - "side_b": same structure for the other document

For each candidate, decide ONE verdict:
- "AGREE": both statements express compatible or mutually reinforcing
  positions.
- "UNRELATED": the statements talk about different things; the
  structural signal is coincidental.
- "TENSION": the statements diverge in strength, scope, or framing
  without directly contradicting each other (e.g., one mandates what
  the other merely recommends; definitions emphasise different aspects).
- "CONTRADICT": the statements cannot both be satisfied at the same time.

Rules:
- Judge ONLY from the provided texts. Do not invent context.
- Standards rarely contradict outright; prefer TENSION over CONTRADICT
  unless the texts are truly incompatible.
- The rationale must be ONE sentence, naming both documents, written
  for a human reviewer.

You MUST respond with:
{
  "verdicts": [
    {"id": <string>, "verdict": <verdict>, "rationale": <string>},
    ...
  ]
}

Hard requirements:
- Output strict JSON only. No markdown.
- Exactly one verdict per input candidate with the original id.

Now judge these candidates:
$candidates_json
\end{lstlisting}

% ----------------------------------------------------------
\subsection{Stage 2b — Keyword Synonym Expansion}
\label{app:prompt-synonyms}

Used by Stage~2b to expand a user-supplied focus keyword into up to
ten synonyms for the KeyBERT relevance gate.

\begin{lstlisting}[language={},breaklines=true,caption={},numbers=none,xleftmargin=0pt]
Return only strict JSON. Do not explain or reason.

Generate up to 10 close synonyms, alternatives, and natural
morphological variants for this search keyword. These terms will be
used for keyword expansion in a semantic search filter.

Required JSON shape:
{"synonyms": ["...", "...", ...]}

Rules:
- Keep the list to a maximum of 10 elements.
- Include close morphological or derivational variants when
  natural/common.
- Do not include broad unrelated concepts.
- Output JSON only.

Search keyword:
$keyword
\end{lstlisting}

% ────────────────────────────────────────────────────────────
\section{Predicate Vocabulary}
\label{app:predicates}

Table~\ref{tab:predicates} lists the complete controlled predicate
vocabulary used in Stage~5 triplet extraction, organised by the four
ontology families.  Every predicate extracted by the LLM must be an
exact string match against this list; off-vocabulary predicates are
tagged \texttt{Needs schema review} and held for human override.

\begin{table}[h]
\centering
\small
\caption{Controlled predicate vocabulary (Stage~5).}
\label{tab:predicates}
\begin{tabular}{ll}
\toprule
\textbf{Family} & \textbf{Predicates} \\
\midrule
Normative
  & \texttt{Requires}, \texttt{Recommends}, \texttt{Permits},\\
  & \texttt{Prohibits}, \texttt{Defines}, \texttt{Provides},\\
  & \texttt{AppliesTo}, \texttt{Constrains}, \texttt{Measures} \\[4pt]
Trustworthy AI
  & \texttt{IsDimensionOf}, \texttt{ContributesTo},\\
  & \texttt{MightMitigate}, \texttt{MightIntroduce},\\
  & \texttt{SurfacesRisk}, \texttt{IsThreatTo} \\[4pt]
Data Science
  & \texttt{Implements}, \texttt{HasParameter}, \texttt{HasInput},\\
  & \texttt{HasOutput}, \texttt{Evaluates}, \texttt{Performs} \\[4pt]
Taxonomy
  & \texttt{IsSubclassOf}, \texttt{IsEquivalentTo}, \texttt{IsA} \\
\bottomrule
\end{tabular}
\end{table}

% ────────────────────────────────────────────────────────────
\section{Five-Strategy Verifier Specification}
\label{app:verifier}

Table~\ref{tab:verifier} gives the pass/fail thresholds and checking
logic for each of the five post-upsert verification strategies
(Stage~8).  Each strategy runs in an isolated exception handler so a
single failure degrades to a per-strategy verdict rather than aborting
the full audit.

\begin{table}[h]
\centering
\small
\caption{Stage~8 verification strategy specifications.}
\label{tab:verifier}
\begin{tabularx}{\columnwidth}{llX}
\toprule
\textbf{ID} & \textbf{Threshold} & \textbf{Check} \\
\midrule
S1 Coverage      & $\tau \geq 0.90$ & $\geq$90\% of source quality sentences are reflected in at least one triple, verified via provenance record matching. \\[4pt]
S2 Correctness   & $\tau = 1.0$     & All nodes have non-placeholder names; all edges use a predicate from the declared vocabulary. \\[4pt]
S3 Consistency   & $\tau = 1.0$     & No $(S, O)$ pair is simultaneously linked by both an affirmative predicate (\texttt{Requires}, \texttt{Recommends}, \texttt{Permits}) and a prohibitive one (\texttt{Prohibits}). \\[4pt]
S4 Completeness  & $\tau = 1.0$     & Every triple has a non-empty subject, object, and predicate, and at least one provenance record. \\[4pt]
S5 Minimality    & $\tau = 1.0$     & No duplicate $(S, P, O)$ edge exists in the graph. \\
\bottomrule
\end{tabularx}
\end{table}

% ────────────────────────────────────────────────────────────
\section{Pre-Merge KB Preview Categories}
\label{app:preview-categories}

Before any graph write, each proposed triple is assigned one of five
pre-merge categories (Stage~6a):

\begin{description}
  \item[\textsc{new}] The $(S, P, O)$ combination is not yet in the
    graph.  Concepts may already exist as nodes, but the relation is new.
  \item[\textsc{existing}] An identical or equivalent triple is already
    present with the same predicate and compatible deontic modality.
  \item[\textsc{tension}] The same $(S, P, O)$ exists but with a
    different obligation strength (e.g.\ this document says
    \textsc{mandatory}, the graph says \textsc{recommended}).
  \item[\textsc{conflict}] The same subject–object pair is already linked
    by a directly contradicting predicate (e.g.\ \texttt{Requires}
    vs.\ \texttt{Prohibits}).
  \item[\textsc{related}] The same concept pair appears in the graph but
    via a different, non-contradicting predicate.
\end{description}

\textsc{conflict} and \textsc{tension} items are sorted to the top of
the review queue.  Confirmed conflicts are materialised as
\texttt{:Conflict} nodes in the graph for downstream compliance
reasoning.

% ────────────────────────────────────────────────────────────
\section{Comparison Engine: Conflict Type Definitions}
\label{app:conflict-types}

The multi-document comparison engine detects four structural conflict
types before invoking the LLM judge (Appendix~\ref{app:prompt-conflict-judge}):

\begin{description}
  \item[Obligation conflict] The same $(S, P, O)$ triple appears with
    different deontic modalities across two documents
    (e.g.\ \textsc{mandatory} in one, \textsc{recommended} in another).
  \item[Definition divergence] The same concept is the target of a
    \texttt{Defines} edge in two documents, but the defining sentences
    differ semantically above a configurable threshold.
  \item[Taxonomy conflict] A subsumption reversal: document~A asserts
    $X$ \texttt{IsSubclassOf} $Y$ while document~B asserts $Y$
    \texttt{IsSubclassOf} $X$.
  \item[Value conflict] The same $(S, P, \cdot)$ pattern appears in
    both documents but with different object values for a
    value-bearing predicate.
\end{description}

% ────────────────────────────────────────────────────────────
\section{Few-Shot Example Counts per Prompt}
\label{app:fewshot-counts}

Each prompt template is injected with a fixed set of hand-authored
few-shot examples at runtime.  Table~\ref{tab:fewshot} documents the
example counts.  All examples are domain-agnostic (not drawn from the
EU AI Act, ISO/IEC~42001, or ISO/IEC~5259-3 to avoid data leakage).

\begin{table}[h]
\centering
\small
\caption{Few-shot example counts per prompt.}
\label{tab:fewshot}
\begin{tabular}{lc}
\toprule
\textbf{Prompt} & \textbf{Examples} \\
\midrule
Stage 2c — Atomic Decomposition (batch)       & 3 \\
Stage 4  — Novelty Classification (batch)     & 4 \\
Stage 5  — Triplet Extraction (batch)         & 4 \\
Stage 5  — Triplet Extraction (neutral/MINE)  & 3 \\
Stage 2b — Keyword Synonym Expansion          & 0 \\
Stage 6  — Conflict Judge                     & 0 \\
\bottomrule
\end{tabular}
\end{table}
